\begin{frame}[t]{\ev{\tagM}Six findings from two experiments}
\vspace{.1cm}
{\renewcommand{\arraystretch}{1.45}\setlength{\tabcolsep}{4pt}
\begin{tabularx}{\textwidth}{@{}>{\color{Teal}\bfseries}l >{\bfseries}l Y@{}}
\toprule
1 & {Sandbox setup took little time} & setup was under 1\% of a 43-minute work order\\
2 & {A locked test folder did not limit the repair} & a repair rewrote code outside the plan\\
3 & {Passing checks can test nothing} & CI skipped every step; a review approved untested code\\
4 & {One model performed all four steps} & one model wrote, reviewed, fixed and approved\\
5 & {Teardown can destroy the result} & an approved patch was lost with its sandbox\\
6 & {No fork was available} & the loop requested a fork; none was offered\\
\bottomrule
\end{tabularx}\par}
\claim{Across seven surfaces, trusted components should run outside the sandbox.}
\sources{Our experiments, Sep 2026: two work orders in Docker Sandboxes, models served through Databricks. Records: \href{https://github.com/zozo123/ariflow-swfactory}{github.com/zozo123/ariflow-swfactory}.}
\note{[0:45] The runs gave six findings. Sandbox setup took under one percent of work order time. A locked test folder did not limit the repair. Passing checks can test nothing. One model in all four steps gives one review. Teardown can destroy an approved result. No provider offered a fork. Across the seven surfaces, trusted components should run outside the sandbox. Next we consider copying the sandbox.}
\end{frame}
